\section{One Command, and What It Talks To}
\label{sec:ch09-one-command}
\index{composition!running it|(}%

Everything in this chapter comes out of one command:

\begin{minted}{text}
npx wgc router compose -i federation/mosaic.yaml -o federation/supergraph.json
\end{minted}

\index{composition!wgc router compose@\code{wgc router compose}}%
\code{federation/mosaic.yaml} was written in
chapter~\ref{ch:the-first-cut-extracting-catalog}, which composed with it on
every verification run and printed none of the result. It names the two
subgraphs, says where each one answers, and points at the file holding each
one's schema:

\begin{minted}{yaml}
version: 1

subgraphs:
  - name: catalog
    routing_url: http://localhost:5101/graphql
    schema:
      file: ../schema/catalog.graphql

  - name: mosaic
    routing_url: http://localhost:5100/graphql
    schema:
      file: ../schema/mosaic.graphql
\end{minted}

That is the entire input. Two schema files and a note of which URL each one
belongs to.

\index{composition!offline}%
The command needs no account, no registry and no network. Cosmo's documentation
says it \enquote{does not interact with the control plane and completely runs
locally}~\autocite{cosmo2026routercompose}, and with schema files as the input
it does not contact the subgraphs either. The routing URLs go into the output
for the router to use later. The composer never visits them, and
chapter~\ref{ch:how-a-federated-query-actually-runs} printed the one request it
does make when you give it the other kind of input.

\index{composition!subgraphs need not be running}%
Which has a consequence worth stating plainly, because it is the difference
between composition being a build step and being an integration test: neither
service has to be running. Both were stopped while I wrote this section. A
composition failure is a statement about two text files.

\subsection*{Where the schemas come from, and what that costs}

\index{composition!file input versus introspection}%
The \code{file:} form is a choice with a failure mode attached. The composer
reads what is committed under \code{schema/}, and what is committed is what each
service published through \code{\_service} on the day somebody exported it.
Change a subgraph's C\# without re-exporting, and the composer will go on
composing a graph that no longer exists.

\code{introspection:} in place of \code{schema:} points it at a live endpoint
instead, and chapter~\ref{ch:how-a-federated-query-actually-runs} printed the
request it sends. That form cannot drift, and it cannot run in a build that has
no services in it either.

Mosaic keeps the files and closes the gap a different way, which is
section~\ref{sec:ch09-composition-as-a-gate}: the verification script asks each
running service for its schema and fails if what it publishes is not the
committed file. The composer reads a file, and something else makes sure the
file is true. Chapter~\ref{ch:ci-cd-for-a-federated-graph} is where that
arrangement becomes a pipeline with a registry in it, and where the drift check
stops being local.

\subsection*{What the command answers with}

Success is quiet:

\begin{minted}[fontsize=\footnotesize,breaklines]{text}
Router execution config successfully written to "federation/supergraph.json".
\end{minted}

Failure is a table, and the exit code is 1. The tables in
sections~\ref{sec:ch09-satisfiability} and~\ref{sec:ch09-reading-the-errors} are
drawn with box-drawing characters and hard-wrapped to a fixed width, so every
error in this chapter is printed as its message rather than as its box. The
words are the composer's; the line breaks around them are the command-line
tool's, and I have taken them out.

\index{composition!determinism}%
One property of the command matters more than it sounds like it should:
composition is deterministic. I composed the same pair twice and compared the
two outputs byte for byte with \code{cmp}, and they are identical. That is what
makes the rest of this chapter possible. An artefact you can commit is an
artefact a build can diff, and section~\ref{sec:ch09-composition-as-a-gate} does
exactly that.

Now the file itself, which is not the file I expected.

\index{composition!running it|)}%
